\chapter{Six-Dimensional Superconformal Theories}
\label{ch:susy-six-dimensional-scft}

\ConventionsMixed

Six-dimensional superconformal theories are QFTs whose central examples do
not admit ordinary Lagrangian descriptions at the fixed point.  They therefore
force a separation between the definition of a theory and the use of
effective field variables on a branch of vacua.  This chapter records the
objects that can be defined without assuming a microscopic Lagrangian:
superconformal representation data, tensor-branch effective actions, anomaly
polynomials, BPS strings, compactification tests, and extended operators.

\section{Six-Dimensional Supersymmetry Data}

Let the Lorentz group be \(\operatorname{Spin}(1,5)\).  A chiral spinor has four complex
components before reality conditions.  Six-dimensional supersymmetry is
labeled by \((\mathcal N_+,\mathcal N_-)\), the numbers of positive and
negative chirality symplectic-Majorana-Weyl supercharges.  The cases most
central for local interacting fixed points are \((1,0)\) and \((2,0)\).

For \((1,0)\), the R-symmetry is \(SU(2)_R\).  For \((2,0)\), the
R-symmetry is \(USp(4)_R\).  Superconformal symmetry enlarges the bosonic
group to
\[
  SO(2,6)\times G_R ,
\]
with \(G_R=SU(2)_R\) or \(USp(4)_R\).  Local operators are organized as
positive-energy representations of the corresponding superconformal algebra.
This statement concerns Hilbert-space representation theory on the cylinder;
it does not supply off-shell fields.

\section{Multiplets And Branch Coordinates}

A \((1,0)\) tensor multiplet contains a two-form gauge potential \(B\) with
self-dual field strength, a real scalar \(\phi\), and fermions.  A
hypermultiplet contains scalars in quaternionic representations and fermions.
A vector multiplet can appear in gauge-theory descriptions on a tensor branch,
but it is not a conformal gauge field at the origin of a six-dimensional
interacting fixed point.

On a tensor branch, expectation values of tensor-multiplet scalars define
coordinates
\[
  \phi^I\in\mathcal C ,
\]
where \(\mathcal C\) is a cone modulo any discrete identifications.  The
low-energy action contains
\[
  \frac12 G_{IJ}\,d\phi^I\wedge *d\phi^J
  +
  \frac12 G_{IJ}\,H^I\wedge *H^J
  +
  B^I\wedge X_I^{(4)}
  +\cdots ,
\]
with self-duality imposed as a constraint on the equations of motion or by a
chosen self-dual-field formalism.  The four-form \(X_I^{(4)}\) encodes gauge
and gravitational couplings of the branch EFT.

\section{Anomaly Polynomial}

For a six-dimensional theory coupled to background metric and R-symmetry
fields, the anomaly is encoded by an eight-form
\[
  I_8
  \in
  H^8(BG_{\mathrm{background}},\mathbb Q)
\]
whose descent gives anomalous variations of the generating functional.  For a
\((1,0)\) theory one may write
\[
  I_8
  =
  \alpha\,c_2(R)^2
  +
  \beta\,c_2(R)p_1(T)
  +
  \gamma\,p_1(T)^2
  +
  \delta\,p_2(T)
  +
  I_8^{\mathrm{flavor}},
\]
where \(c_2(R)\) is the second Chern class of the \(SU(2)_R\) bundle and
\(p_i(T)\) are Pontryagin classes of the tangent bundle.  The coefficients
are part of the QFT datum.

On a tensor branch, anomaly matching takes the form
\[
  I_8^{\mathrm{UV}}-I_8^{\mathrm{branch}}
  =
  \frac12\Omega^{IJ}X_I^{(4)}X_J^{(4)} ,
\]
where \(\Omega^{IJ}\) is the inverse tensor pairing.  The right-hand side is
produced by Green--Schwarz couplings \(B^I\wedge X_I^{(4)}\).  This equation
is a computable constraint on branch EFTs and on proposed fixed points.

\section{BPS Strings}

Tensor multiplets couple electrically to strings.  A string with charge
\[
  q\in\Lambda_{\mathrm{string}}
\]
has tension
\[
  T_q=q_I\phi^I
\]
on the tensor branch, in a chamber where this pairing is positive.  At the
superconformal point some tensions vanish, and the corresponding string
degrees of freedom are not particles in six dimensions.

The worldsheet theory of a BPS string carries anomalies determined by inflow
from the bulk Green--Schwarz term.  If the string worldsheet is \(\Sigma\),
the anomaly polynomial has the form
\[
  I_4(\Sigma,q)
  =
  q_I X_I^{(4)}
  +
  \frac12 q_I\Omega^{IJ}q_J\,\chi(N_\Sigma)
  +\cdots ,
\]
where \(N_\Sigma\) is the normal bundle.  The omitted terms must be fixed by
the chosen R-symmetry and tangent-bundle conventions, not by analogy.

\section{Compactification Tests}

Compactification on a smooth compact manifold \(Y\) gives lower-dimensional
QFTs whose protected data test the six-dimensional proposal.  For a circle
of radius \(R\), a \((2,0)\) theory of type \(\mathfrak g\) yields five
dimensional maximally supersymmetric Yang--Mills on its Coulomb branch with
coupling scaling as
\[
  g_5^2\propto R .
\]
The proportionality depends on trace and instanton-charge conventions.  The
Kaluza--Klein momentum along \(S^1\) is identified with instanton number in
the five-dimensional description.

For compactification on a Riemann surface \(C\), partial twists use an
embedding of the spin connection of \(C\) into \(G_R\).  The resulting
four-dimensional theory must have matching anomalies, flavor symmetries,
defect descendants, and moduli-space dimensions.  These are tests of a
six-dimensional construction, not definitions of the parent theory by
themselves.

\section{Status Ledger}

A six-dimensional SCFT datum in this monograph consists of:
\[
\begin{gathered}
  \text{superconformal representation category},\\
  \text{tensor-branch EFT with self-dual fields and Green--Schwarz terms},\\
  \text{anomaly polynomial and flavor-background data},\\
  \text{BPS string charge lattice and inflow},\\
  \text{compactification tests and extended-operator descendants}.
\end{gathered}
\]
No item in this list is dispensable.  A proposed fixed point that supplies
only a branch EFT is incomplete; a proposed non-Lagrangian object that
supplies only protected indices is also incomplete.

\begin{openproblem}[Six-dimensional construction]
\label{op:six-dimensional-construction}
Construct an interacting six-dimensional SCFT as a local QFT with Hilbert
space, local operator algebra, stress tensor, anomaly polynomial, tensor
branch, BPS strings, and compactification functoriality in one framework.
\end{openproblem}
